\documentclass[11pt,a4paper]{article}
\usepackage[utf8]{inputenc}
\usepackage{amsmath,amssymb,amsfonts,amsthm}
\usepackage{geometry}
\geometry{margin=1in}
\usepackage{hyperref}
\usepackage{cite}
\usepackage{bm}
\usepackage{verbatim}

\title{\textbf{A Major-Minor Mean Field Game Framework for Crowd Navigation and Dynamic Goal Evasion}}
\date{\today}

\begin{document}

\maketitle

\begin{abstract}
This document establishes a formal mathematical framework for modeling the interaction between a continuum of minor players (a crowd) and dynamic major players (evading or capacity-constrained goals) within a bounded spatial domain containing obstacles. The system is formulated as a Major-Minor Mean Field Game (PE-MFG). We detail the backward Hamilton-Jacobi-Bellman (HJB) equation governing optimal crowd navigation, the forward Kolmogorov-Fokker-Planck (KFP) equation modeling crowd mass transport, the repulsive potential force field governing dynamic evading goals, and the centralized goal capacity saturation scheme. Finally, a damped Picard fixed-point scheme with post-saturation trajectory clamping is presented alongside the corresponding software architecture and code implementation details.
\end{abstract}

\section{Introduction \& Domain Configuration}
Let $\Omega \subset \mathbb{R}^2$ be a bounded spatial domain with boundary $\partial\Omega = \Gamma_{\text{wall}} \cup \Gamma_{\text{door}}(t)$, where $\Gamma_{\text{wall}}$ represents solid obstacles enforcing zero-flux Neumann conditions and $\Gamma_{\text{door}}(t)$ denotes time-dependent exit boundaries associated with active goals. The temporal horizon is $t \in [0, T]$.

The system consists of two agent classes:
\begin{enumerate}
    \item \textbf{Minor Players (Crowd):} A continuum of statistically identical agents whose macroscopic spatial density at time $t$ and location $x \in \Omega$ is denoted by $m(t, x) \ge 0$.
    \item \textbf{Major Players (Goals):} A discrete collection of $G$ goal objects $\{y_g(t)\}_{g=1}^G \subset \Omega$. Goals can be stationary landing pads, prescribed parametric trajectories, or reactive evading goals that actively flee from crowd density. Goals may also possess finite mass capacity ceilings $C_{\text{max}, g}$.
\end{enumerate}

\section{Minor Player (Crowd) Dynamics}

Each crowd member minimizes a cost functional penalizing spatial distance from active goals, spatial congestion, and control effort.

\subsection{Hamilton-Jacobi-Bellman (HJB) Equation}
Let $u(t, x)$ denote the value function (expected cost-to-go) of a crowd member at position $x$ and time $t$. The value function satisfies the backward-in-time HJB equation:
\begin{equation}
\label{eq:hjb}
-\frac{\partial u}{\partial t} - \nu \Delta u + H(x, m, \nabla u) = f(x, t), \quad (t, x) \in (0, T) \times \Omega,
\end{equation}
subject to the terminal condition $u(T, x) = +f(x, T)$ and Dirichlet boundary condition $u(t, x) = 0$ for $x \in \Gamma_{\text{door}}(t)$.

The congestion-dependent Hamiltonian $H(x, m, p)$ models crowd mobility reduction at higher densities:
\begin{equation}
H(x, m, p) = -\frac{\alpha}{(1 + m)^\beta} |p|^2 + h_0,
\end{equation}
where $\alpha, \beta > 0$ control mobility reduction and $h_0$ is a baseline running cost.

\subsection{Active Goal Selection \& Running Cost}
Let $\mathcal{A}(t) \subseteq \{1, \dots, G\}$ denote the set of goals that remain active and unsaturated at time $t$. The running cost $f(x, t)$ penalizes distance to the nearest unsaturated goal and incorporates an optional repulsive penalty $S(x, t)$ around saturated goals:
\begin{equation}
\label{eq:running_cost}
f(x, t) = \chi \min_{g \in \mathcal{A}(t)} \|x - y_g(t)\|^2 + S(x, t),
\end{equation}
where $\chi > 0$ is the attraction weight. If a goal $g$ is fully saturated ($g \notin \mathcal{A}(t)$), it is removed from the minimum-distance calculation. If a non-zero saturated goal penalty coefficient $w_{\text{sat}} > 0$ is configured, saturated goals exert a local Gaussian repulsive barrier:
\begin{equation}
S(x, t) = w_{\text{sat}} \sum_{g \notin \mathcal{A}(t)} \exp\left( -\frac{\|x - y_g(t)\|^2}{2 \sigma_{\text{sat}}^2} \right).
\end{equation}

\subsection{Kolmogorov-Fokker-Planck (KFP) Equation}
The aggregate crowd density $m(t, x)$ evolves forward in time according to the optimal velocity drift field $\mathbf{v}(t, x) = -\nabla_p H(x, m, \nabla u)$:
\begin{equation}
\label{eq:kfp}
\frac{\partial m}{\partial t} - \nu \Delta m - \text{div}\left( m \, \frac{2 \alpha \nabla u}{(1 + m)^\beta} \right) = 0, \quad (t, x) \in (0, T) \times \Omega,
\end{equation}
subject to the initial density $m(0, x) = m_0(x)$, homogeneous Neumann conditions on $\Gamma_{\text{wall}}$, and Dirichlet absorption ($m(t, x) = 0$) on active exit regions $\Gamma_{\text{door}}(t)$.

\section{Goal Dynamics \& Capacity Saturation}

\subsection{Centralized Saturation State Tracking}
Any goal $g$ configured with a finite capacity limit $C_{\text{max}, g} < \infty$ automatically acts as a mass-absorbing exit ($is\_exit = \text{true}$). The cumulative crowd mass that has exited through the goal region $\Omega_g(t)$ up to time $t$ is given by:
\begin{equation}
C_g(t) = \int_0^t \int_{\Omega_g(s)} m(s, x) \, dx \, ds.
\end{equation}
The saturation step $t_{\text{sat}, g}$ is defined as the earliest time at which cumulative exited mass reaches capacity:
\begin{equation}
t_{\text{sat}, g} = \min \left\{ t \in [0, T] : C_g(t) \ge C_{\text{max}, g} \right\}.
\end{equation}
For all $t \ge t_{\text{sat}, g}$:
\begin{enumerate}
    \item Goal $g$ is marked as saturated ($g \notin \mathcal{A}(t)$) and removed from the HJB attraction field.
    \item The corresponding exit boundary region closes: $\Gamma_{\text{door}, g}(t) = \emptyset$.
    \item The goal trajectory freezes permanently at its saturation coordinates: $y_g(t) = y_g(t_{\text{sat}, g})$.
\end{enumerate}

\subsection{Evasive Goal Trajectories}
For active ($t < t_{\text{sat}, g}$) evading goals, the motion vector $y_g(t)$ is driven by a continuous repulsive force field $F(y_g(t))$ computed from the local crowd density distribution:
\begin{equation}
\label{eq:repulsion_force}
F(y_g(t)) = \int_{\Omega} \frac{y_g(t) - x}{\|y_g(t) - x\|^2 + \epsilon} m(t, x) \, dx,
\end{equation}
where $\epsilon > 0$ is a spatial regularization parameter preventing singularities. The velocity vector is constrained by a maximum speed parameter $v_{\text{max}, g}$:
\begin{equation}
v_g(t) = \begin{cases} 
v_{\text{max}, g} \dfrac{F(y_g(t))}{\|F(y_g(t))\|}, & \text{if } \|F(y_g(t))\| > \delta, \\[8pt]
0, & \text{otherwise.}
\end{cases}
\end{equation}
The spatial position updates forward in time via Euler integration and is projected onto the walkable domain $\Omega \setminus \Omega_{\text{obs}}$ to prevent obstacle collision:
\begin{equation}
y_g(t + \Delta t) = \begin{cases} 
\mathcal{P}_{\Omega \setminus \Omega_{\text{obs}}} \left( y_g(t) + v_g(t) \Delta t \right), & \text{if } t < t_{\text{sat}, g}, \\[6pt]
y_g(t_{\text{sat}, g}), & \text{if } t \ge t_{\text{sat}, g}.
\end{cases}
\end{equation}

\section{Algorithmic Resolution via Extended Picard Iteration}

The coupled HJB-KFP system and dynamic goal trajectories are solved iteratively using a damped Picard fixed-point scheme with under-relaxation parameter $\theta \in (0, 1]$. To prevent trajectory bleeding across iterations, post-saturation positions are explicitly clamped.

\begin{enumerate}
    \item \textbf{Initialization:} Initialize $U^{(0)}(t, x)$, $M^{(0)}(t, x)$, and goal trajectories $Y^{(0)}(t) = \{y_g^{(0)}(t)\}_{g=1}^G$.
    \item \textbf{Step 1 (Goal Motion \& Saturation Update):} Given crowd density $M^{(k-1)}$, compute temporary goal trajectories $Y_{\text{temp}}$ and evaluate cumulative mass absorption $C_g(t)$. Identify saturation steps $t_{\text{sat}, g}$ and assemble the exit door mask $\Gamma_{\text{door}}(t)$.
    \item \textbf{Step 2 (HJB Backward Solve):} Solve Eq.~\eqref{eq:hjb} backward from $t = T$ to $t = 0$ using $M^{(k-1)}$, active goal set $\mathcal{A}(t)$, and $Y_{\text{temp}}$ to obtain $\tilde{U}^{(k)}$. Update value function:
    \begin{equation}
    U^{(k)} = \theta \tilde{U}^{(k)} + (1-\theta) U^{(k-1)}.
    \end{equation}
    \item \textbf{Step 3 (KFP Forward Solve):} Solve Eq.~\eqref{eq:kfp} forward from $t = 0$ to $t = T$ using $U^{(k)}$ and $\Gamma_{\text{door}}(t)$ to obtain $\tilde{M}^{(k)}$. Update crowd density:
    \begin{equation}
    M^{(k)} = \theta \tilde{M}^{(k)} + (1-\theta) M^{(k-1)}.
    \end{equation}
    \item \textbf{Step 4 (Trajectory Relaxation \& Post-Saturation Clamping):} Apply under-relaxation to goal trajectories:
    \begin{equation}
    Y^{(k)} = \theta Y_{\text{temp}} + (1-\theta) Y^{(k-1)}.
    \end{equation}
    For all goals $g$ that reach capacity at $t_{\text{sat}, g} \le T$, enforce hard position clamping for all $t \ge t_{\text{sat}, g}$:
    \begin{equation}
    y_g^{(k)}(t) = y_g(t_{\text{sat}, g}), \quad \forall t \ge t_{\text{sat}, g}.
    \end{equation}
    \item \textbf{Convergence Check:} Repeat until space-time L2 residuals satisfy $\|U^{(k)} - U^{(k-1)}\|_2 < \epsilon$, $\|M^{(k)} - M^{(k-1)}\|_2 < \epsilon$, and $\|Y^{(k)} - Y^{(k-1)}\|_2 < \epsilon$.
\end{enumerate}

\section{Software Architecture \& Code Implementation}

The mathematical formulation is implemented as a modular Python library (\texttt{mfgames}). High-performance finite difference operators are accelerated using Numba JIT compilation, and sparse linear systems are assembled in COO format and solved via Scipy sparse solver routines.

\subsection{Single Source of Truth for Goal Saturation (\texttt{mfgames/objectives.py})}
The \texttt{Goal} class manages goal initialization, trajectory updates, and capacity tracking. In \texttt{Goal.update_positions()}, cumulative mass $C_g(t)$ exiting through each goal region is integrated. Once $C_g(t) \ge C_{\text{max}, g}$, the state array \texttt{is_saturated[k, g]} is set to \texttt{True}, and positions are permanently frozen:

\begin{verbatim}
# Track cumulative mass exiting at step k
if is_exit and np.isfinite(cap):
    if k > 0 and self.is_saturated[k - 1, g]:
        self.is_saturated[k:, g] = True
        if k < self.Nt:
            new_trajectories[k + 1:, g] = [curr_x, curr_y]
        continue

    region = (np.abs(X_grid - curr_x) <= Dx) & (np.abs(Y_grid - curr_y) <= Dy)
    mass_exiting_at_k = np.sum(M_k[region]) * Dx * Dy
    cumulative_mass[g] += mass_exiting_at_k

    if cumulative_mass[g] >= cap:
        self.is_saturated[k:, g] = True
        if self.saturation_steps[g] > k:
            self.saturation_steps[g] = k
        if k < self.Nt:
            new_trajectories[k + 1:, g] = [curr_x, curr_y]
        continue
\end{verbatim}

\subsection{Active Goal Filtering \& Trajectory Clamping (\texttt{mfgames/problem.py})}
The \texttt{MFGSolver} class coordinates the coupled HJB-KFP solver loop. In \texttt{compute_running_cost()}, the solver queries \texttt{Goal.is_saturated} directly to filter out saturated goals from the active attraction set $\mathcal{A}(t)$ and apply optional repulsive penalties:

\begin{verbatim}
def compute_running_cost(self, goal_positions_k, k=None):
    X, Y = self.pde_mesh.X, self.pde_mesh.Y
    active_distances = []
    saturated_repulsion = np.zeros((self.Nx, self.Ny))

    for g_idx, (gx, gy) in enumerate(goal_positions_k):
        is_active = True
        if k is not None and self.goal is not None and self.goal.has_capacity_limits:
            if self.goal.is_saturated[k, g_idx]:
                is_active = False

        if is_active:
            dist_sq = (X - gx) ** 2 + (Y - gy) ** 2
            active_distances.append(dist_sq)
        elif self.saturated_goal_penalty > 0.0:
            dist_sq = (X - gx) ** 2 + (Y - gy) ** 2
            sigma_sq = (10.0 * self.Dx) ** 2
            saturated_repulsion += self.saturated_goal_penalty * np.exp(-dist_sq / (2.0 * sigma_sq))

    if active_distances:
        min_dist_sq = np.minimum.reduce(active_distances)
    else:
        min_dist_sq = np.zeros((self.Nx, self.Ny))

    return self.running_cost_weight * min_dist_sq + saturated_repulsion
\end{verbatim}

During the macro Picard iteration loop in \texttt{run_picard_system()}, under-relaxation is applied to goal trajectories, followed by an explicit clamp to freeze positions for all $t \ge t_{\text{sat}, g}$:

\begin{verbatim}
# Step 5: Relax goal trajectories and clamp post-saturation positions across all k >= sat_k
if self.goal is not None and Y_temp is not None:
    Y_new = self.thetaUM * Y_temp + (1.0 - self.thetaUM) * self.goal.Y_trajectories

    for g_idx in range(self.goal.num_goals):
        sat_k = self.goal.saturation_steps[g_idx]
        if sat_k <= self.Nt:
            Y_new[sat_k:, g_idx, :] = Y_temp[sat_k, g_idx, :]

    self.goal.Y_trajectories = np.copy(Y_new)
\end{verbatim}

\subsection{Visualizing Saturation States (\texttt{mfgames/plotting.py})}
The \texttt{MFGPlotter} class handles snapshot dashboards and video generation. In \texttt{_draw_goals()}, it reads \texttt{Goal.is_saturated[t_idx, g_idx]} directly to color-code markers dynamically:

\begin{verbatim}
def _draw_goals(self, ax, t_idx=0):
    if self.evader_trajectories is not None and self.goal_instance is not None:
        for g_idx, g_info in enumerate(self.goal_instance.goals):
            pos = self.evader_trajectories[t_idx, g_idx]
            is_sat = self.goal_instance.is_saturated[t_idx, g_idx]
            marker_color = '#ff2222' if is_sat else '#00f2fe'

            ax.scatter(pos[0], pos[1], color=marker_color, marker='X', s=80, edgecolor='black', zorder=10)
\end{verbatim}

\begin{thebibliography}{99}

\bibitem{caines2017mean}
P.~E. Caines, M.~Huang, and R.~P. Malham{\'e}, 
\emph{Mean field games}, 
Handbook of Dynamic Game Theory, Springer, pp.~345--372, 2017.

\bibitem{carmona2016probabilistic}
R.~Carmona and X.~Zhu, 
\emph{A probabilistic approach to mean field games with major and minor players}, 
The Annals of Applied Probability, vol.~26, no.~3, pp.~1535--1580, 2016.

\bibitem{huang2016backward}
J.~Huang, S.~Wang, and Z.~Wu, 
\emph{Backward-forward linear-quadratic mean-field games with major and minor agents}, 
Probability, Uncertainty and Quantitative Risk, vol.~1, no.~1, p.~8, 2016.

\bibitem{bergault2024mean}
P.~Bergault, P.~Cardaliaguet, and C.~Rainer, 
\emph{Mean field games in a Stackelberg problem with an informed major player}, 
SIAM Journal on Control and Optimization, vol.~62, no.~3, pp.~1737--1765, 2024.

\bibitem{badakaya2022game}
A.~J. Badakaya, A.~S. Halliru, and J.~Adamu, 
\emph{Game value for a pursuit-evasion differential game problem in a Hilbert space}, 
Journal of Dynamics \& Games, vol.~9, no.~1, pp.~1--15, 2022.

\bibitem{wan2021improved}
K.~Wan, D.~Wu, Y.~Zhai, B.~Li, X.~Gao, and Z.~Hu, 
\emph{An improved approach towards multi-agent pursuit--evasion game decision-making using deep reinforcement learning}, 
Entropy, vol.~23, no.~11, p.~1433, 2021.

\end{thebibliography}

\end{document}